\section{The actual large-gap discrete series and the real-stage join}
\label{sec:discrete-series-core}

This section constructs the discrete series needed by the later row parameters.
It does not assert the source's universal \texttt{DiscreteSeriesExist} or
\texttt{FormalDegreeFormula\(_2\)}.  The construction below uses finite compact
representation theory, integration and Hilbert-space arguments.  The degree
calculation in DC18--DC22 has one additional, explicitly unreconstructed
classical premise: Harish--Chandra's formal-degree polynomial theorem, stated
precisely in DC18.  Its Haar constant is calculated here.  The volume part of
the final join retains GV's named Tamagawa theorem.  No statement below
constructs an entire \texttt{Core} from a volume or a representation alone.

Number coordinates by $0,\ldots,n-1$.  Write
\[
 G_\epsilon=\{g:g^*J_\epsilon g=J_\epsilon\},\qquad
 K_\epsilon=G_\epsilon\cap U(n),\qquad
 p=\#\epsilon^{-1}(+),\quad q=\#\epsilon^{-1}(-),\quad q_0=\min(p,q).
\]
The positive roots are $e_a-e_b$ with $a<b$.  Separate them into compact and
noncompact roots according as $\epsilon_a=\epsilon_b$ or not, and let
$\rho_c,\rho_n$ be their half-sums.  Thus
$\rho=\rho_c+\rho_n$, $\rho_i=(n-1)/2-i$.  Inner products on Hilbert spaces
are conjugate linear in the first argument.  A parameter is \emph{regular
integral} here precisely when
\[
 \lambda_0>\cdots>\lambda_{n-1},\qquad
 \lambda_i-(n-1)/2\in\mathbb Z.                                  \tag{DC.1}
\]
Put $\Lambda=\lambda+\rho_n-\rho_c$,
$c_\lambda=\|\lambda\|^2-\|\rho\|^2$ and
$c_K(\nu)=\|\nu+\rho_c\|^2-\|\rho_c\|^2$.  A noncompact gap $h$ means
$|\lambda_a-\lambda_b|\ge h$ whenever $\epsilon_a\ne\epsilon_b$.
For the constructive existence assertions we require
\[
 h>H_\epsilon:=\max\{npq+1,\ 16nq_0+8\sqrt{2q_0}\}.             \tag{DC.2}
\]
The smaller bound $h>npq+1$ suffices for the rigidity assertions once a
model has been constructed.  Compact signatures are treated separately;
there is no division by $q_0$ in that case.

\paragraph{DC1: the exact compact type and its integral construction.}
\label{claim:DC1}
The weight $\Lambda$ is integral and dominant on each sign block.  Indeed
$\lambda+\rho$ is integral and
$\Lambda=\lambda+\rho-2\rho_c$.  If $a<b$ are consecutive indices in one
sign block, and $r$ opposite-sign indices lie between them, then
\[
 \Lambda_a-\Lambda_b=\lambda_a-\lambda_b+r-1\ge0.
\]
The last inequality follows from (DC.1): consecutive distinct entries of
$\lambda$ differ by a positive integer.  For an integral dominant weight
$\alpha$ of $U(d)$, subtract its last entry to get a partition.  The Young
symmetrizer applied to the corresponding tensor power of $\mathbb C^d$,
followed by the determinant power $\det^{\alpha_{d-1}}$, gives the compact
irreducible representation of highest weight $\alpha$.  Averaging its
positive inner product over $U(d)$ makes it unitary.  The usual elementary
highest-weight argument proves irreducibility: a highest vector is a scalar
multiple of the row-symmetrized, column-alternating tableau vector, and
lowering matrix units generate every tableau.  Apply this independently
to the two blocks.  We obtain a finite-dimensional unitary $K_\epsilon$-module
$T$, with a unit highest vector $w$, the exact object called \texttt{KType}
in the source.  This uses no noncompact representation-existence theorem.

\paragraph{DC2: the principal-minor function and the normalized average.}
\label{claim:DC2}
Set $\mu_i=\lambda_i+\rho_i\in\mathbb Z$, put $\mu_n=0$, and let
$\Delta_j(z)$ be the leading $j$ by $j$ minor.  On the open set where these
minors are nonzero define
\[
 \Phi_\mu(z)=\prod_{j=1}^n\Delta_j(z)^{-(\mu_{j-1}-\mu_j)},
 \qquad
 P(z)=d_T\int_K\Phi_\mu(k^{-1}zk)\,T(k)P_wT(k)^{-1}\,dk,         \tag{DC.3}
\]
where $d_T=\dim T$, $dk$ is probability Haar and $P_w v=w\langle w,v\rangle$.
Every leading minor of a positive definite matrix is positive.  Compactness
of $K$ therefore gives an open complex neighborhood of the positive cone
on which the second integral is holomorphic; all its derivatives may be
taken under the integral on a smaller neighborhood.  Triangular
multiplication gives
\[
 \Phi_\mu(lzu)=\prod_i(l_{ii}u_{ii})^{-\mu_i}\Phi_\mu(z)
                                                                  \tag{DC.4}
\]
for lower triangular $l$ and upper triangular $u$.  Schur averaging gives
$d_T\int_K T(k)P_wT(k)^{-1}dk=1$, hence $P(1)=1$.
For $X=X^*$ with $JX=-XJ$, one has $e^{-X}=Je^XJ$; diagonal conjugation
by $J$ fixes every leading minor.  Consequently $P(e^{-X})=P(e^X)$.
The source formula with an inverse matrix in the argument of $\Phi$ is
this identity, not a reciprocal of $\Phi$.

\paragraph{DC3: the compact differential identities and the global kernel.}
\label{claim:DC3}
For every block-diagonal complex matrix $Y$, the average satisfies
\[
 dP_z(zY)=-dT(Y)P(z),\qquad dP_z(Yz)=-P(z)dT(Y).                \tag{DC.5}
\]
Here is the integration-by-parts calculation that supplies the essential
weight shift.  Write $B=k^{-1}Yk$, and decompose $B$ into its strictly lower,
diagonal and strictly upper parts inside the two blocks.  Differentiating
(DC.4) gives the diagonal contribution $-2\sum_i\mu_i B_{ii}$ to the
variation in $Yz+zY$.  Moving the lower-minus-upper variation onto
$T(k)P_wT(k)^{-1}$ differentiates the two representation factors.
The integral of a compact left-invariant derivative is zero.  Differentiating
its coefficient $B_{ab}$ contributes
\[
 \sum_{a\ne b,\,\epsilon_a=\epsilon_b}
       \operatorname{sgn}(a-b)(B_{aa}-B_{bb})
       =-4\sum_a(\rho_c)_aB_{aa}.
\]
The remaining diagonal action is exactly $-2\sum_i\Lambda_iB_{ii}$,
because $\mu=\Lambda+2\rho_c$.  Thus
$dP_z(Yz+zY)=-dT(Y)P(z)-P(z)dT(Y)$ for Hermitian block $Y$.
Differentiating conjugation equivariance supplies the difference identity;
adding and subtracting gives (DC.5) for skew-Hermitian block $Y$.
Complex linearity and holomorphic continuation from the real positive
cone give the asserted complex identities.

The polar decomposition $g=k e^X$, $k\in K$, $X=X^*$, $JX=-XJ$, is unique.
Define
\[
             \Psi(g)=T(k)P(e^X).                               \tag{DC.6}
\]
Then $\Psi(k_1gk_2)=T(k_1)\Psi(g)T(k_2)$ and $\Psi(1)=1$.
For a fixed $g_0=k_0e^{X_0}$, a local logarithm of $k(g)^{-1}k_0$ and
(DC.5) show that, for $g$ sufficiently close to $g_0$,
\[
                  \Psi(g)=T(k_0)P(g^{-1}k_0).                  \tag{DC.7}
\]
Explicitly, the derivative of $T(e^{tY})P(ze^{tY})$ is zero.
Integrate this identity along the short compact logarithm path and use
$P(e^{-X})=P(e^X)$.  The right side of (DC.7) is a holomorphic function
of an ambient invertible complex matrix near $g_0$.  This proves the
literal local ambient smoothness required by the source.

\paragraph{DC4: central differential operators have the actual parameter.}
\label{claim:DC4}
Let $E_{ab}$ denote the complexified right derivatives in the source
convention.  The kernel has
\[
 \sum_{a,b}E_{ab}E_{ba}\Psi=c_\lambda\Psi,\qquad
 \mathcal C\bigl(Y+(n-1)/2\bigr)\Psi
       =\prod_i(Y+\lambda_i)\Psi.                              \tag{DC.8}
\]
In the second formula the literal column-determinant convention is
\[
 \mathcal C(u)=\sum_{\sigma\in S_n}\operatorname{sgn}(\sigma)
   \prod_{j=0}^{n-1}\bigl(E_{\sigma(j),j}
                      +\delta_{\sigma(j),j}(u-j)\bigr),
\]
with factors composed in increasing $j$ from left to right (the
rightmost factor acts first).  Thus in rank one it is $Y+E_{00}$;
for the trivial compact representation in rank two, $\lambda=\rho$
and $u=Y+1/2$, it is $(Y+1/2)(Y-1/2)$.
These are actual scalar differential identities, not a classification
assumption.  To check them, use (DC.7), differentiate under the compact
integral, and move conjugations through the invariant contracted matrix
units.  On $\Phi_\mu$, strictly lower left derivatives vanish and diagonal
left derivatives are $-\mu_i$ by (DC.4).  Commuting an upper derivative
past a lower one uses
$[E_{ab},E_{cd}]=\delta_{bc}E_{ad}-\delta_{ad}E_{cb}$.
For the quadratic sum the diagonal result is
$\sum_i\mu_i^2-2\sum_i\rho_i\mu_i
=\|\lambda\|^2-\|\rho\|^2$ after the inversion in (DC.7).
For the determinant, expand by columns and commute each strictly lower
factor to the right.  Every off-diagonal commutator cancels the term of
the permutation differing by that transposition; its diagonal commutator
in column $i$ contributes the shift $n-1-i$.  Only the diagonal product
survives.  Substituting $\mu_i=\lambda_i+(n-1)/2-i$ and reversing left
derivatives through inversion yields the displayed product.  The same
calculation with the first $r$ columns proves the identity inductively,
so it is valid in every rank.  These coefficients generate the center:
their highest PBW symbols are the coefficients of the ordinary
characteristic polynomial, which generate conjugation-invariant
polynomials on matrices; subtracting that leading symbol and inducting
on the filtration proves the assertion.  In particular the full infinitesimal
parameter of this kernel is the unordered list $\{\lambda_i\}$; the
quadratic eigenvalue alone would not establish that fact.

\paragraph{DC5: a logarithmic minor estimate with explicit constants.}
\label{claim:DC5}
If $w_j\ge0$, $\sum w_j=1$, $\sum w_jy_j=0$ and $|y_j|\le R$, then
\[
 \log\sum_jw_je^{y_j}\ge
       \tfrac14\min(1,R^{-1})\sum_jw_jy_j^2.                   \tag{DC.9}
\]
Take $c=\min(1,R^{-1})$, use $e^u\ge1+u+u^2/3$ for $|u|\le1$,
then $\log(1+t)\ge3t/4$ for $0\le t\le1/3$.  Jensen's inequality
for the concave power $x^c$ passes from $cy_j$ to $y_j$ and proves
(DC.9).  If $R=0$ both sides vanish.

Let $Y=k^{-1}Xk$ with $X$ as in DC3 and put $r^2=\operatorname{tr}(X^2)$.
Diagonalize $Y=U\operatorname{diag}(y_i)U^*$.  Cauchy--Binet gives
\[
 \Delta_m(e^Y)=\sum_{|I|=m}|\det U_{[m],I}|^2
                              \exp\bigl(\sum_{i\in I}y_i\bigr).
\]
The coefficients sum to one.  The weighted mean is
$\operatorname{tr}(Y_{[m],[m]})=0$; its variance is
$\sum_{b<m\le a}|Y_{ab}|^2$.  The mean and variance identities follow by
expanding respectively the first and second derivatives at zero of
$\det((e^{tY})_{[m],[m]})$; the second derivative is
$\operatorname{tr}((Y^2)_{[m],[m]})-\operatorname{tr}(Y_{[m],[m]}^2)$.
The rank of $Y$ is at most $2q_0$, so
$|\sum_{i\in I}y_i|\le\sqrt{2q_0}\,r$ by Cauchy--Schwarz.  Hence
\[
 \log\Delta_m(e^Y)\ge\tfrac14
  \min\{1,(\sqrt{2q_0}r)^{-1}\}\sum_{b<m\le a}|Y_{ab}|^2.       \tag{DC.10}
\]
At $X=0$ this is read by continuity.

\paragraph{DC6: decay, square integrability and the exact large-gap bound.}
\label{claim:DC6}
For $q_0>0$, (DC.10) telescopes in the exponents of (DC.3).
The determinant term is one because $\operatorname{tr}X=0$.  Moreover
\[
 \sum_{m<n}(\mu_{m-1}-\mu_m)
                 \sum_{b<m\le a}|Y_{ab}|^2
 =\sum_{b<a}(\mu_b-\mu_a)|Y_{ab}|^2\ge\tfrac h2 r^2.
\]
Only opposite-sign entries of $Y$ are nonzero.  Since the integrand
projectors are positive and their normalized average is the identity,
there is no extraneous $\dim T$ in the operator bound:
\[
 \|\Psi(ke^X)\|\le
 \exp\left[-\tfrac h8\min\{r^2,r/\sqrt{2q_0}\}\right].        \tag{DC.11}
\]
For completeness the polar Jacobian is the determinant of
$\sinh(\operatorname{ad}X)/\operatorname{ad}X$ on the noncompact tangent
space.  Differentiating $k e^X$ proves this formula, including its
continuous value at zero.  Diagonalizing $X$, whose nonzero eigenvalues
occur in opposite pairs, and using $\sinh(t)/t\le e^{|t|}$ bounds it by
$\exp(n\sqrt{2q_0}\,r)$.  Polar Lebesgue measure is ordinary measure
on $\mathbb C^{pq}$; its shell factor is polynomial in $r$.
Thus (DC.11) is integrable for $h>16nq_0$, is square integrable for
$h>8nq_0$, and its product with $e^{sr}$ is square integrable when
$h>8nq_0+8s\sqrt{2q_0}$ (a sufficient bound).  In particular (DC.2)
supplies all unweighted integrability needed here.  Any other Haar
measure is a positive scalar multiple and has the same finiteness
properties.

\paragraph{DC7: the kernel entries are genuinely smooth vectors in $L^2$.}
\label{claim:DC7}
Fix $v_0\ne0$ in $T$ and define
$j(v)(g)=\langle v_0,\Psi(g)v\rangle$.
The finite-dimensional space $j(T)$ consists of smooth $L^2$ functions,
is stable under the compact right derivatives, and has Casimir
$c_\lambda$ by DC3--DC6.  These facts imply that \emph{every} right
derivative word of $j(v)$ lies in $L^2$.

Here is the required global regularity argument.  Choose orthonormal real
compact and noncompact Lie bases $Y_i,X_j$ for the trace form.  On smooth
functions the positive elliptic operator is
$L=-\sum_iY_i^2-\sum_jX_j^2=2C_K-C_G$, with the conventional compact
$C_K=-\sum_iY_i^2$.  Use compactly supported cutoffs $\chi_R$ on the
complete invariant metric, equal to one on the radius-$R$ ball, with
$|d\chi_R|\le C/R$.  Integration by parts for $\chi_R f$, followed by
Cauchy--Schwarz and absorption of half the derivative term, bounds
$\sum\|\chi_R X_jf\|_2^2$ by a fixed multiple of
$\sum\|Y_if\|_2^2+(1+|c_\lambda|)\|f\|_2^2$ and a vanishing cutoff
error.  No derivative integrability was assumed in this compactly
supported calculation.  Fatou gives the first noncompact derivatives
in $L^2$.  The space of words of length at most $m$ applied to $j(T)$
is finite-dimensional and compact-stable, and the central Casimir
still has eigenvalue $c_\lambda$.  Repeating the same argument proves
the assertion by induction on $m$.
Finally the fundamental theorem of calculus along a one-parameter
subgroup, now in $L^2$, identifies these derivatives with the generators
of the right regular representation.  Iterating proves smoothness of
the orbit map.  Evaluation at the identity and continuity show $j$ is
nonzero.  Its kernel is compact-invariant, so compact irreducibility
makes $j$ injective.

\paragraph{DC8: the finite tensor calculation that eliminates lower types.}
\label{claim:DC8}
For a unitary infinitesimal representation the matrix units satisfy
$E_{ab}^*=E_{ba}$ on compact roots and $E_{ab}^*=-E_{ba}$ on
noncompact roots.  If $u$ is a compact highest vector of weight $\nu$
and $C_Gu=c_\lambda u$, expansion of the Casimir gives
\[
 (c_\lambda+\|\rho_c-\rho_n\|^2)\|u\|^2
 =\|\nu+\rho_c-\rho_n\|^2\|u\|^2
       -2\sum_{a<b,\,\epsilon_a\ne\epsilon_b}\|E_{ba}u\|^2.    \tag{DC.12}
\]
The proof pairs the $(a,b)$ and $(b,a)$ summands, uses the adjoint signs,
and replaces their commutator by $E_{aa}-E_{bb}$.  In particular no
compact highest vector of weight $\Lambda-(e_a-e_b)$ occurs if
$\lambda_a-\lambda_b>1+2\langle\rho_c,\rho_n\rangle$.
Counting ordered sign pairs gives
$2\langle\rho_c,\rho_n\rangle\le npq$, so the first inequality in
(DC.2) suffices.

There is a useful finite-dimensional version of this assertion.  On
$\mathfrak p_{\mathbb C}\otimes T$ put
\[
 A=\sum_{\epsilon_a=\epsilon_b}
                \operatorname{ad}(E_{ab})\otimes dT(E_{ba}).
\]
The operator is self-adjoint.  The elementary standard/dual tensor rule
in each compact block shows that every summand has highest weight
$\Lambda+e_r-e_s$, $\epsilon_r\ne\epsilon_s$, when that weight is
dominant.  This rule follows by locating a largest nonzero coordinate
of a highest vector in $\mathbb C^d\otimes V$ and solving the upper
matrix-unit equations; the dual rule is identical with indices reversed.
Since the compact Casimir on $\mathfrak p_{\mathbb C}$ is $n$, its
eigenvalue on this summand is
\[
 \theta_{rs}=\tfrac12\{c_K(\Lambda+e_r-e_s)-c_K(\Lambda)-n\}
 =\lambda_r+(\rho_n)_r-\lambda_s-(\rho_n)_s+1-n/2.
\]
For $r<s$ it is at least $h-n+1>0$; for $r>s$ it is at most $-h-1<0$.
Let $W_-,W_+$ be the negative and positive spectral subspaces.  The
multiplication map $X\otimes v\mapsto E(X)v$ kills $W_-$ by (DC.12).
This finite set of first-order relations is the Schmid relation used
below; its vanishing has been proved, not assumed.

\paragraph{DC9: coercive PBW layers and the minimum-norm condition.}
\label{claim:DC9}
Let $F_j$ be the span of words of at most $j$ noncompact derivatives
applied to $T$.  PBW and compact stability show that these exhaust the
generated Lie module.  If a nonzero highest vector of weight $\nu$
occurs in $F_j/F_{j-1}$, $j\ge1$, then
\[
 c_K(\nu)\ge c_K(\Lambda)+2j(h-n+1).                           \tag{DC.13}
\]
This conclusion requires only the first-order relation $W_-\mapsto0$.
Indeed choose the minimal-norm lift in
$\mathfrak p_{\mathbb C}^{\otimes j}\otimes T$ orthogonal to the
preimage of $F_{j-1}$.  Moving a negative slot to the last position
introduces commutators of smaller word length, so every negative-slot
tensor maps into $F_{j-1}$.  Every slice of the chosen lift is therefore
in $W_+$.  Expanding the total compact Casimir gives
\[
 C_K^{\mathrm{total}}=C_K^{\mathrm{tensor}}+c_K(\Lambda)
                                      +2\sum_{i=1}^j A_i.
\]
The first term is nonnegative, and DC8 bounds each $A_i$ below by
$h-n+1$.  The minimal lift is compact-equivariant, proving (DC.13).

In the tensor rule $\nu-\Lambda$ is a sum of $j$ noncompact roots.
Each has absolute pairing at most $(n-2)/2$ with $\rho_c$.
Consequently
\[
 \|\nu+2\rho_c\|^2-\|\Lambda+2\rho_c\|^2
 \ge j(2h-3n+4)>0.                                           \tag{DC.14}
\]
The last inequality follows from $p,q\ge1$, $pq\ge n-1$ and
$h>npq+1$.  Every fixed compact highest weight can occur only in
layers
$j\le[c_K(\nu)-c_K(\Lambda)]/[2(h-n+1)]$.
The original compact type occurs only in $F_0$ and with multiplicity
one.  These are finite-dimensional inequalities; no general
Blattner formula or discrete-series classification has been used.

\paragraph{DC10: closing the Lie module does not introduce extra compact types.}
\label{claim:DC10}
In the right regular $L^2$ space, close the module generated by $j(T)$
and call the resulting Hilbert space $H_\lambda$.  It is invariant under
$G_\epsilon$.  To justify passage from Lie operators to the group, choose
$A\ge0$ bounding the finitely many weights of $T$ and the Casimir scalar.
Every highest weight in $F_j$ has norm at most $A+\sqrt2 j$; the energy
identity used in DC7 therefore implies, for $X\in\mathfrak p$ and
$v\in F_k$,
\[
 \|E(X)^m v\|\le \|X\|_{\rm HS}^{m}
       \prod_{r=0}^{m-1}(A+\sqrt2(k+r))\,\|v\|.               \tag{DC.15}
\]
The Taylor series converges whenever
$|t|\sqrt2\|X\|_{\rm HS}<1$.  The integral remainder estimate, using
unitarity at every base point of the orbit, identifies its sum with
the actual group orbit.  Subdividing a finite interval gives invariance
under $e^X$; compact invariance and polar decomposition give invariance
under $G$.  Invariance under inverses gives equality.

For a compact irreducible type $\nu$, its compact character projector
is a bounded operator.  Apply it to the dense algebraic module.  DC9
places its range in one fixed finite-dimensional PBW layer, which is
closed.  The same is therefore true of its range on $H_\lambda$.
This proves admissibility, smoothness of all compact-finite vectors,
and persistence of the unique copy $j(T)$.  In particular completing
the module has not silently enlarged its lowest-type multiplicity.

\paragraph{DC11: irreducibility and identification of the actual coefficient.}
\label{claim:DC11}
Let $R$ be the restriction of the right regular representation to
$H_\lambda$.  If $S$ is a closed invariant subspace, its orthogonal
projection commutes with $G$ and hence with $K$.  By DC10 it maps
$j(T)$ into itself, and Schur's lemma makes its restriction a scalar.
Since it is a projection, that scalar is zero or one.  The space $j(T)$
generates $H_\lambda$, so $S=0$ or $H_\lambda$.  Thus $R$ is irreducible.

The invariant inner product on $j(T)$ is a positive scalar multiple of
that of $T$.  If $w$ is unit highest, put $a=\|j(w)\|_2^2>0$.
The exact coefficient is
\[
           \langle j(v'),R(g)j(v)\rangle
                  =a\langle v',\Psi(g)v\rangle.               \tag{DC.16}
\]
Here is a direct identification, avoiding an assumed positive-definite
kernel.  Project $R(g)j(v)$ onto $j(T)$, say to $j(u)$, and let the
orthogonal remainder be $f$.  Its compact matrix-coefficient averages
give compact maps $T\to H_\lambda$ orthogonal to $j(T)$; by uniqueness
they are zero.  The functions being averaged have continuous
representatives.  Their zero $L^2$ classes are therefore zero pointwise,
since Haar measure is positive on nonempty open sets.  Evaluate at $1$,
using $\Psi(kg)=T(k)\Psi(g)$ and $\Psi(k)=T(k)$.  The relevant compact
Schur integral is
$\|v_0\|^2\|\Psi(g)v-u\|^2/d_T$, so $u=\Psi(g)v$.
The inner product with $j(v')$ proves (DC.16).

\paragraph{DC12: all coefficients, Schur orthogonality and the source degree.}
\label{claim:DC12}
Equation (DC.16) and DC6 give a nonzero square-integrable coefficient
of $R$.  This implies every coefficient is square integrable and proves
the full four-vector orthogonality identity.  The functional-analytic
step can be seen explicitly.  For fixed $v$ define the closed operator
\[
 T_vu(g)=\langle v,R(g)u\rangle
\]
on the vectors for which this function lies in $L^2(G)$.  For an initial
nonzero square-integrable coefficient its domain contains the dense
span of translates of a nonzero vector.  Closedness follows by an
$L^2$ subsequence and pointwise continuity of the coefficient in $u$.
The positive self-adjoint operator $T_v^*T_v$ commutes with $R$.
All its spectral projections do too, so irreducibility makes it a
finite nonzero scalar times the identity.  Hence $T_v$ is bounded
and defined on the entire Hilbert space.  Inversion and complex
conjugation exchange the two coefficient slots; repeating the same
argument and polarizing in both slots gives
\[
 \int_G\langle w,R(g)v\rangle
       \overline{\langle w',R(g)v'\rangle}\,dg
       ={\langle v',v\rangle\langle w,w'\rangle\over d_R},
       \qquad d_R>0.                                         \tag{DC.17}
\]
Products are integrable by Cauchy--Schwarz.  This is precisely
\texttt{IsFdeg}.  Setting all four vectors equal to a nonzero $w$
shows that the source quotient \texttt{fdegOf} equals $d_R$ and is
independent of $w$.  Replacing $dg$ by $c\,dg$, $c>0$, replaces
$d_R$ by $d_R/c$.  DC9 gives the minimum Vogan norm required in the
literal \texttt{IsDS}; completeness and nontriviality are already
established.  We have constructed that entire predicate under (DC.2).

\paragraph{DC13: rigidity for the source's lowest-type definition.}
\label{claim:DC13}
An arbitrary representation satisfying the source's \texttt{IsDS}
for the same $(\epsilon,\lambda)$ is unitarily equivalent to the
constructed model under (DC.2).  This assertion is needed because
\texttt{IsDS} specifies a minimum compact type rather than a full
infinitesimal character.

First smooth a nonzero copy of the prescribed compact type by a
compact-conjugation-invariant approximate identity.  On that finite
copy the convolution is arbitrarily close to the identity, hence
injective, compact-equivariant and smooth.  It supplies a smooth
copy $V$ of $T$.  A first-layer type $\Lambda-e_a+e_b$, $a<b$,
would have
\[
 \|\Lambda-e_a+e_b+2\rho_c\|^2
       -\|\Lambda+2\rho_c\|^2
       =2-2(\lambda_a-\lambda_b+b-a)<0,
\]
contrary to the literal minimum-norm condition.  Thus the same
negative tensor space $W_-$ maps to zero.  The Schmid version of DC9
applies even before a Casimir scalar is known.

To recover that scalar without a classification theorem, take the
diagonal copy of $T$ in the direct sum of $V$ and the constructed
model's initial type.  Both projections kill $W_-$, so DC9 implies
that the generated diagonal module contains just one highest line
of weight $\Lambda$.  The central quadratic Casimir maps that line
to itself.  Projection to the constructed model determines its
scalar to be $c_\lambda$; projection to $V$ gives the same scalar
there.  Centrality extends it to the whole generated module.
DC10 closes this diagonal module to an invariant Hilbert subspace,
and DC11 makes it irreducible.  Its two coordinate projections are
nonzero intertwiners with dense range onto the two irreducible
representations.  The positive operators obtained by composing a
projection with its adjoint are scalar by Schur's lemma, hence each
projection is a positive multiple of a unitary equivalence.
Their ratio is the desired unitary intertwiner.  Consequently all
the central eigenvalues in (DC.8), as well as the formal degree for
any fixed Haar measure, hold for every literal \texttt{IsDS} at
these large-gap parameters.  This proof uses neither a general
lowest-type classification nor a universal existence theorem.

\paragraph{DC14: the later row weights satisfy the gap in the literal order.}
\label{claim:DC14}
The actual row rank is $n=R.r=k+m-1$.  Its ordered parameter consists
of $v_0,\ldots,v_{k-2}$, followed by $-b_0,\ldots,-b_{m-1}$;
the head has sign $+$ and the tail has alternating signs starting
with $+$.  Interlacing gives $v_{k-2}>u_{k-1}>-b_0$.
Therefore every opposite-sign head/tail difference is larger than
$b_1-b_0$, and every opposite-sign tail difference is at least the
least successive $b$-gap.  Head/head pairs are compact.
The source parity assumptions make $k,m$ even, $n$ odd and every
entry of this row parameter integral, so (DC.1) holds.  Strict
decrease follows from the same interlacing and increasing $b$'s.

The context has $b_j=b_0+jh_0$.  Its actual eventual quantifier
supplies, for arbitrary fixed $A\ge0,N,M$, simultaneously
\[
                   M\le m,\qquad A n^N\le h_0.                \tag{DC.18}
\]
For clarity, the source witnesses may be taken as $\epsilon_0=1$,
zero transported alternating data, $c_{D,0}=0$, $c_{h,0}=6N+1$ and
$x_0=\max(\max(4,A),M^4)$.  Its rank estimates give $n\le x^6$
and $m\ge x^{1/4}$, while $h_0\ge x^{c_h}$.  Thus these witnesses
prove (DC.18) in the original nested order, uniformly for every
later coefficient field and every same-predata object.
Take $A=16,N=19,M=1801$.  Then $n\ge1801$ and
$h_0\ge16n^{19}>H_\epsilon$, since $p,q,q_0\le n$.
This is the literal \texttt{GapHyp} route, and it preserves $B$, its
host, $R$, every $\mathrm{hc}_\tau$ and the supplied stage objects.

\paragraph{DC15: actual row existence, including the endpoint clauses.}
\label{claim:DC15}
For every $D$ satisfying \texttt{GapHyp}, and every real place $\tau$,
DC1--DC14 give a complete nonzero Hilbert representation
\[
             \operatorname{IsDS}(\epsilon_{SH},
                                  \lambda_{SH}(R,\mathrm{hc}_\tau),\sigma_\tau).
                                                                  \tag{DC.19}
\]
The choices are over the same finite set of real places; they do not
alter $D$.  The construction also gives finite compact isotypic
spaces, smooth compact-finite vectors, a unique lowest compact type,
the exact central parameter, and uniqueness among the source's
\texttt{IsDS} representations at these parameters.  For compact
signatures the representation of DC1 itself has these properties;
ordinary compact Schur orthogonality gives its degree.  In ranks
zero and one the empty products and the one-dimensional characters
give degree one for probability Haar.  Thus no rank-zero or
compact-signature exceptional branch is hidden in (DC.19).

\paragraph{DC16: the precise compact-dual quotient density.}
\label{claim:DC16}
Write $W_d=\prod_{j=0}^{d-1}j!$, with $W_0=1$.  GV5 proved that
compact probability Haar in ordinary skew-Hermitian coordinates
has density $1/v_d$, where $v_d=2^d\pi^{d(d+1)/2}/W_d$.
The source's compact-dual convention carries this same tangent
density to $U(p,q)$.  In block coordinates the compact factor has
probability density $1/(v_pv_q)$, so quotient measure on
$G/K$ has origin density
\[
             C_{pq}={v_pv_q\over v_n}
                    =\pi^{-pq}{W_n\over W_pW_q}.              \tag{DC.20}
\]
Identify $G/K$ with
$\mathcal D_{pq}=\{Z\in M_{p,q}(\mathbb C):1-Z^*Z>0\}$.
The derivative at the origin sends the off-diagonal Lie block
to $Z$ with Jacobian one.  The derivative of
$Z\mapsto(aZ+b)(cZ+d)^{-1}$ and the identity
\[
 1-(gZ)^*(gZ)=(cZ+d)^{-*}(1-Z^*Z)(cZ+d)^{-1}
\]
give the invariant quotient density
$C_{pq}\det(1-Z^*Z)^{-n}dZ$.
In particular (DC.20) concerns the actual noncompact quotient
measure induced from the compact-dual tangent convention, not
just the total volume of the compact group.

\paragraph{DC17: a scalar holomorphic model pins the normalization.}
\label{claim:DC17}
For an integer $s>n-1$, use holomorphic functions on $\mathcal D_{pq}$
with norm proportional to
\[
          \int_{\mathcal D_{pq}}|f(Z)|^2
                            \det(1-Z^*Z)^{s-n}\,dZ.           \tag{DC.21}
\]
The action
$f(Z)\mapsto\det(c_{g^{-1}}Z+d_{g^{-1}})^{-s}f(g^{-1}Z)$
is unitary by the preceding determinant and Jacobian identities.
Polynomials are dense: dilations $f(rZ)$ converge in the norm,
and their Taylor series converge on the smaller closed domains.
The central circle of $K$ separates homogeneous degrees; the
constant type occurs only in degree zero.  Translates of the
constant span densely.  Indeed averaging on that circle computes
the value at zero, and after a group translation computes the
value at any point; a function orthogonal to all such translates
vanishes everywhere.  The unique constant compact type and this
cyclicity prove irreducibility by the projection argument of DC11.

The beta integral in (DC.21) is elementary.  Remove the last row
$z$ of $Z$ and put $A=1-Z'^*Z'$.  The substitution $z=uA^{1/2}$
has real Jacobian $\det A$, and
$\det(1-Z^*Z)=\det A(1-\|u\|^2)$.  Radial integration in
$\mathbb C^q$ yields
\[
 I_{p,q}(a):=\int_{\mathcal D_{pq}}\det(1-Z^*Z)^a\,dZ
 =\pi^{pq}\prod_{r=0}^{p-1}{\Gamma(a+r+1)\over\Gamma(a+q+r+1)},
 \quad a>-1.                                                  \tag{DC.22}
\]
This follows by $p$ iterations of
$I_{p,q}(a)=\pi^q\Gamma(a+1)I_{p-1,q}(a+1)/\Gamma(a+q+1)$.
For a unit constant vector the squared coefficient is
$\det(1-Z^*Z)^s$.  Equations (DC.20)--(DC.22) and DC12 therefore give
\[
 d_s={1\over C_{pq}I_{p,q}(s-n)}
       ={W_pW_q\over W_n}
          \prod_{i=0}^{p-1}\prod_{j=0}^{q-1}(s-n+1+i+j).       \tag{DC.23}
\]
The central parameter can be checked on the constant vector by
triangular matrix units.  Order the negative block before the
positive block.  The two lists are
\[
 s+(q-p-1)/2-j\quad(0\le j<q),\qquad
 (n-1)/2-i\quad(0\le i<p).
\]
They are strictly ordered for $s>n-1$, have the required parity,
and their within-block Vandermondes are $W_q,W_p$.
Their cross differences have product equal to the double product
in (DC.23).  Hence this explicit model has
$d_s=|\prod_{a<b}(\lambda_a-\lambda_b)|/W_n$ for exactly the
compact-dual Haar measure.  No formal-degree theorem is used to
compute this calibration example.

\paragraph{DC18: the named degree theorem and its fully determined constant.}
\label{claim:DC18}
\emph{Explicit unreconstructed classical input HC-Deg.}
For a connected semisimple real group with finite center and a
compact Cartan, the formal degrees of its discrete series are a
positive group-and-Haar constant times the absolute product of
positive-root pairings of their regular infinitesimal parameters;
the constant is independent of the parameter and of the packet.
We use Harish--Chandra, \emph{Discrete series for semisimple Lie
groups. II: Explicit determination of the characters}, Acta
Mathematica 116 (1966), 1--111, \S40 Corollary 2 (p.~94) and
\S41 Theorem 16 (p.~96), DOI \texttt{10.1007/BF02392813}.
The publisher-version copy is available at
\url{https://repository.ias.ac.in/30833/1/336.pdf}.
The statement concerns the regular infinitesimal parameter; a
quadratic Casimir or a positive real number is not a substitute.

The reduction to $U(p,q)$ adds no unknown normalization depending
on the representation.  The finite covering
$U(1)\times SU(p,q)\to U(p,q)$ pulls an irreducible representation
back to a central character tensor an irreducible representation
of $SU(p,q)$.  To see irreducibility of the latter, a subspace
stable under $SU(p,q)$ is stable under the scalar central action
and hence under $U(p,q)$.  Probability measure on $U(1)$ contributes
one to the squared coefficient integral.  Finite-cover Haar
conversion is a fixed scalar independent of the character.
Root pairings depend only on differences of the $\lambda_i$.
HC-Deg therefore gives
$d_\lambda=A_{pq}|\prod_{a<b}(\lambda_a-\lambda_b)|$ for one
positive constant $A_{pq}$.  DC17 computes one such representation
and forces $A_{pq}=1/W_n$.  Compact groups give the same result
directly by the Weyl dimension formula and probability Schur
orthogonality.  Thus, \emph{conditionally on HC-Deg},
\[
               d_\lambda=
            \prod_{a<b}{|\lambda_a-\lambda_b|\over b-a}
                    =\operatorname{fdegHC}(\lambda).          \tag{DC.24}
\]
The equality $\prod_{a<b}(b-a)=W_n$ follows by collecting factors
with fixed $b$.  Every use of (DC.24) below retains HC-Deg.

\paragraph{DC19: the exact degree for every literal large-gap representation.}
\label{claim:DC19}
Assume HC-Deg, (DC.1) and (DC.2).  DC4 identifies the actual
infinitesimal parameter of the constructed model, so DC18 applies
to it.  DC13 transfers the result to \emph{every} representation
satisfying the literal \texttt{IsDS} at this parameter.
If $\mathrm{CdScaleIs}(cd,sc)$ holds, its scaled source measure
is the same compact-dual Haar measure under every allowed sign
sorting, by GV17 and uniqueness of Haar measure.  Consequently,
for each nonzero vector $w$,
\[
 \operatorname{fdegOf}(sc(n,\epsilon)\mu_G,\sigma,w)
             =\operatorname{fdegHC}(\lambda).                 \tag{DC.25}
\]
Conversely any $d$ satisfying \texttt{IsFdeg} for that same
representation and measure equals this number: insert $w$ in
all four slots of (DC.17) and divide by $\|w\|^4>0$.
Thus both the existential and the universal degree clauses
needed at the actual row parameters are supplied.  This is a
restricted supplier with all the source's vector and representation
quantifiers; it is not the all-parameter \texttt{FormalDegreeFormula\(_2\)}.

\paragraph{DC20: all four fields of the actual real-stage record.}
\label{claim:DC20}
Assume HC-Deg.  Let the \emph{same} $D$ satisfy
\texttt{GapHyp}, $\mathrm{VolTieG}(D)$,
\[
 D.s.cdScale=\mathrm{trueScale}(D.cd),\qquad
 D.dsig(\tau)=\mathrm{trueDsig}(D.R,D.\mathrm{hc}_\tau)
                    \quad\hbox{for every }\tau.               \tag{DC.26}
\]
Then $\mathrm{RealTied}(D)$ holds.  The field \texttt{cdScale}
is GV17.  The field \texttt{dsig} is DC15 together with
(DC.24), under the literal measure in (DC.26).  The field
\texttt{dsig\_all} is DC19, with its universal quantifier over
all Hilbert representations and all their asserted degrees.
The field \texttt{vol} is the supplied \texttt{VolTieG}.
This establishes the precise large-gap replacement for
\texttt{LevelThreeToFour.realOf}.  It neither assumes nor proves
the universal \texttt{RealOf} record, whose input $D$ has no
gap restriction.

When $D.vol$ is the actual \texttt{baseVol}/\texttt{trueVol}
on the same host, spaces and $cd$, GV18--GV19 supply the last
field under their explicitly retained Tamagawa-number theorem.
In this use DC20 is conditional on both HC-Deg and that theorem;
the two premises do different work.  An arbitrary stored
positive volume does not satisfy the field merely by positivity.

\paragraph{DC21: the same-stage Core join and the fields still needed.}
\label{claim:DC21}
Fix the literal $K,X,f,T$ of \texttt{coreOfLay}, a layout
$Y:\mathrm{Lay}(X.\mathrm{toD7})$, and its exact
\texttt{Over}, \texttt{RecOK}, \texttt{LedOK} witnesses.
Assume the following remaining inputs on these very objects:
the given $\xi$ and twist set $Wt$, the same-map \texttt{MuIsW},
forced inclusion of $g_0,g_1,d$ and the specified small primes,
\texttt{GatNum}, the untwisted row vanishing, \texttt{MtStableAt},
the primitive type-character power identities, and
$NChoice.Q=T.Q$, $NChoice.k'=T.k$.  Suppose also that $X.\mathrm{toD7}$
satisfies the gap and (DC.26), with actual volume tied as in DC20.
Then the literal \texttt{Core} on $X,f,T$ exists, conditionally
on HC-Deg and on any named theorem used to supply that volume.

Here is the field audit, rather than an implicit reassembly.
The data fields \texttt{xi}, \texttt{Wt} and \texttt{mu} are the
ones supplied; \texttt{isOver}, \texttt{forced}, \texttt{gatNum},
\texttt{mtStable} and \texttt{chiT} are the corresponding assumptions.
The object field is literally $O=\texttt{layObj}(Y)$, with the same
layout $Y$ throughout these assignments.
\texttt{fval} is \texttt{LedOK.1} with \texttt{MuIsW.off};
\texttt{col} is the same layout's prime, embedding, Minkowski and
type-place fields together with \texttt{RecOK};
\texttt{gy} is its ramification/filling fields with \texttt{LedOK.2}.
The nine \texttt{row} fields are the literal row layout, with only
\texttt{untw} supplied separately.  The theta/scalar fields
\texttt{th\_io}, \texttt{gauge\_def}, \texttt{del}, \texttt{etaF}
are read from that same layout and \texttt{RecOK}.
\texttt{real} is DC20; \texttt{dsig\_eq} is (DC.26).
\texttt{nch\_tie} is the layout's actual local-character tie and
\texttt{nch\_Q} is the asserted conductor-support/exponent pair.
This accounts for all twenty-one fields.  In particular it does
not turn a bare empty-twist stage into a geometric core: the
\texttt{forced} field already requires the actual distinguished
prime to occur.  Changing a stored spectral or volume component
requires a new layout or an explicit dependent transport theorem;
this argument never reuses a layout on a different $X$ by notation.

\paragraph{DC22: a restricted level-three-to-four replacement and its frontier.}
\label{claim:DC22}
Assume HC-Deg and retain the literal source hypotheses
\texttt{LevelThreeCondition}, \texttt{SatNew}, \texttt{PackEx},
the actual volume-to-mass implication, and the real idempotent
tie.  For the source's \texttt{SExt.none} branch, the
\texttt{step\_of}/\texttt{lv4\_of} argument works with DC20 in
place of the universal \texttt{RealOf} premise.  The replacement
is justified at \emph{both} uses of that premise.  First DC14
supplies the gap in the original eventual order.  The explicit
\texttt{exists\_sStage\_true} sets the spectral scalar,
\texttt{trueScale} and \texttt{trueDsig}, and supplies the ordinary
local tie on the same $B,C,W$; one uses the stipulated
\texttt{LevelThreeCondition} to move the level-three data to
this new stage.  DC20 supplies its real record.  \texttt{SatNew}
then returns $Z_s$ and its universal \texttt{SameL} comparison
property.  The source constructs the replacement stage using
the actual old divisibility lattices and heights, transports
columns through \texttt{ColRe.exists\_colsTiedSym}, and applies
the returned \texttt{SameL}/carrier isometries.  The row rank and
every $\mathrm{hc}_\tau$ are unchanged, so DC14's gap and
(DC.26) still hold at the second real-record use; the new volume
tie is furnished by the new stage's actual \texttt{GIs2} witness.
DC20 applies again.  This proves the restricted transition with
all remaining hypotheses displayed.

This conclusion removes the two universal discrete-series wrappers
from this \emph{particular} later-stage route.  It does not prove
\texttt{SatNew}: its spectral construction still requires the
actual cusp finite-dimensionality, rapid-decay, infinitesimal
character and uniqueness inputs in their exact scopes, or new
proofs of the restricted variants used there.  It does not prove
\texttt{PackEx}, a new compatible nonempty motive/twist layout,
or unconditional \texttt{ExTrue}.  The all-parameter existence
and degree wrappers, and the source's optional internal
translation/Laplace proof of degree, are retained as distinct
targets.  No claim in this section is a proof of RH or GRH.
